% TOP_PROBLEM: {"id": 20001108, "title": "Order lifting for sparse additive bases of polynomial values", "category_id": 2, "category": {"id": 2, "name": "combinatorics", "display_name": "Combinatorics", "description": "Counting problems, graph theory, discrete structures.", "slug": "combinatorics", "order_index": 2, "created_at": "2026-07-31T15:26:25.671Z"}, "status": "partially_solved", "problem_number": "AIM-COMBINATORICS-0233", "set_id": 14, "statusQualification": "Fixes the natural-number convention as nonnegative and assumes the polynomial sequence lies there; all sufficiently large positive-leading-coefficient values are covered by the same asymptotic question."}
% FIRST_POSTED: 2026-10-01
% ENTRY_KIND: statement_only
% UnsolvedMath ID 20001108; source catalogue code AIM-COMBINATORICS-0233
% !TeX program = lualatex
% Definition and sources only; this file is not an LLM solution attempt.
\documentclass[11pt,a4paper]{article}
\usepackage[margin=25mm]{geometry}
\usepackage{fontspec}
\setmainfont{lmroman10-regular.otf}[BoldFont=lmroman10-bold.otf,
 ItalicFont=lmroman10-italic.otf,BoldItalicFont=lmroman10-bolditalic.otf]
\usepackage{amsmath,amssymb,mathtools}
\usepackage{unicode-math}
\setmathfont{latinmodern-math.otf}
\AtBeginDocument{\let\setminus\smallsetminus}
\usepackage{xurl,hyperref}
\hypersetup{colorlinks=true,urlcolor=blue}
\setcounter{secnumdepth}{0}
\setlength{\parindent}{0pt}
\setlength{\parskip}{5pt}
\setlength{\emergencystretch}{3em}
\begin{document}
\section*{Order lifting for sparse additive bases of polynomial values}\label{20001108}
{\small UnsolvedMath ID 20001108 · AIM-COMBINATORICS-0233 · Combinatorics · AIM Workshop Problem Lists}
\subsection{Definitions and mathematical statement}
Let $f\in\mathbb Z[x]$ have degree $d\geq2$ and positive leading coefficient, and suppose $f(m)\geq0$ for every positive integer $m$.
Fix an integer $h\geq1$.
For $A\subseteq\mathbb Z_{\geq0}$, put
\[
 hA=\{a_1+\cdots+a_h:a_i\in A\},\qquad
 A(N)=|A\cap\{0,1,\ldots,N\}|.
\]
Summands may repeat, and exactly $h$ are used.
Say that $A$ is an additive basis of order $h$ for the polynomial sequence if $f(m)\in hA$ for every integer $m\geq1$.
It need not represent integers outside that sequence.

Determine the strongest universal lower growth bound for $A(N)$ under this hypothesis, with constants allowed to depend on $f$ and $h$.
The elementary counting scale is $N^{1/(hd)}$: there are order $N^{1/d}$ distinct polynomial values at most $N$, while at most $A(N)^h$ ordered sums can use summands at most $N$.
Nonnegativity ensures that any representation of such a value uses only summands at most $N$.
The workshop asks whether the polynomial structure for $d\geq2$ forces a substantially stronger bound than this counting estimate.
A sharp answer should quantify the improvement and exhibit bases approaching it, or explain when the counting scale is optimal.
The polynomial and the order are fixed while $N\to\infty$; changing the polynomial with $N$ is outside the question.
The use of nonnegative naturals is an explicit convention for the source's word ``naturals''.
\subsection{Short English statement}
Must an additive basis for all values of a nonlinear polynomial be denser than the elementary representation-counting bound predicts?
\subsection{Sources}
\begin{itemize}
\item[S1] ULAM.AI and the UnsolvedMath Contributors, supplied problems.json, record 20001108 (AIM-COMBINATORICS-0233), AIM Workshop Problem Lists. Catalogue identity and source transcription; CC BY 4.0.
\url{https://huggingface.co/datasets/ulamai/UnsolvedMath}.
\item[S2] American Institute of Mathematics, Recent trends in additive combinatorics, item 2.8. Original workshop question underlying AIM-COMBINATORICS-0233.
\url{https://aimath.org/WWN/additivecomb/additivecomb.pdf}.
\end{itemize}
\end{document}
